\uuid{tpeR}
\titre{Application linéaire/Polynôme}
\auteur{hanani ; mzari}
\organisation{exo7}
\licence{CC BY-NC-SA 4.0}
\datecreate{2020-01-20}
\niveau{L1}
\module{Algèbre}
\chapitre{Application linéaire}
\sousChapitre{Définition}
\difficulte{3}
\theme{Applications linéaires, Application linéaire, Polynôme}
\exo7qcmid{lille-2.applications-lineaires.010}

\contenu{
\begin{question}
\qclassification{107.01}
\qtags{motcle=Application linéaire/Polynôme}
On note $\Rr_n[X]$ l'espace des polynômes à coefficients réels de degré $\le n$,  $n\in \Nn$. On considère les deux applications suivantes :
$$\begin{array}{rccc}f:&\Rr_3[X]&\to&\Rr_2[X]\\
& P&\to &R \end{array} \quad \mbox{et} \quad
 \begin{array}{rccc}g:&\Rr_3[X]&\to& \Rr_2[X]\\
& P&\to &Q,\end{array}$$
où $R$ (resp. $Q$) est le reste (resp. le quotient) de la division euclidienne de $P$ par $X^3+1$. Quelles sont les assertions vraies ?
\begin{answers}
\good{$f(0)=0$}.
\good{$f$ est une application linéaire}.
\good{$g(0)=0$}.
\bad{$g$ n'est pas une application linéaire}.
\end{answers}
\begin{explanations} On vérifie que $f$ et $g$ sont linéaires.
\end{explanations}
\end{question}
}
